\section{Cooper's theorem}
 In this section, we derive the Wold decomposition picture of a  $1$-parameter semigroup of isometries in the continuous case. Recall that, in the discrete case, we have already derived Wold decomposition in Thm. \ref{Wold}. Let $\bbr_{+}:=[0,\infty)$. Note that $\bbr_{+}$ is a semigroup. 
 Let us start with a few definitions. 
  \begin{dfn}
 Let $\clh$ be a Hilbert space and let $V:[0,\infty) \to B(\clh)$ be a map. We denote the image of an element $t \in [0,\infty)$ under $V$ by $V_t$. We say that $V:=\{V_t\}_{t \in \bbr_{+}}$ is a strongly continuous semigroup of isometries if 
 \begin{enumerate}
 \item[(1)] for every $t \in \bbr_{+}$, $V_t$ is an isometry,
 \item[(2)] for $s,t \in \bbr_{+}$, $V_{s+t}=V_s V_t$, and
 \item[(3)] for $\xi \in \clh$, the map $\bbr_{+} \ni t \to V_t \xi \in \clh$ is continuous. 
 \end{enumerate}
  \end{dfn}
  Sometime, we call a strongly continuous $1$-parameter semigroup of isometries as an isometric representation of $\bbr_{+}$. Let $V:=\{V_t\}_{t \geq 0}$ be an isometric representation of $\bbr_{+}$ on  a Hilbert space $\clh$. Then, $V$ is said to be \emph{pure} if $V_{t}^{*} \to 0$ strongly, i.e. for $\xi \in \clh$, $V_{t}^{*}\xi \to 0$ as $t \to \infty$. If each $V_t$ is  unitary, then we say that $V$ is a unitary representation of $\bbr_{+}$. 
  
  \begin{xrcs}
  Let $U:=\{U_t\}_{t \geq 0}$ be a unitary representation  of $\bbr_{+}$ on a Hilbert space $\clh$. Then, $U$ can be extended to a unitary representation of the group $\bbr$. 
    \end{xrcs}
  \textit{Hint:} Simply set $U_{-t}:=U_{t}^{*}$ for $t>0$. 
  
  
 \begin{ppsn}
Let $V:=\{V_t\}_{t \geq 0}$ be an isometric representation of $\bbr_{+}$ on a Hilbert space $\clh$. For $t \geq 0$, let $E_t$ be the range projection of $V_t$, i.e. $E_t:=V_{t}V_{t}^{*}$. The following are equivalent. 
\begin{enumerate}
\item[(1)] The isometric representation $V$ is pure. 
\item[(2)] For $\xi \in \clh$, $E_{t}\xi \to 0$ as $t \to \infty$. 
\item[(3)] The intersection $\displaystyle \bigcap_{t \geq 0}V_t\clh=\{0\}$. 
\end{enumerate}
\end{ppsn}  
\textit{Proof.} Let $t \in \bbr_{+}$ be given. Observe that for $\xi \in \clh$, $||E_t\xi|| =||V_tV_{t}^{*}\xi|| \leq ||V_t^{*}\xi||$. Also, for $\xi \in \clh$, 
$
||V_{t}^{*}\xi||=||V_{t}^{*}V_{t}V_{t}^{*}\xi|| \leq ||E_t\xi||$. The equivalence between $(1)$ and $(2)$ is now clear. 

Let $\clh_{\infty}:=\bigcap_{t \geq 0}V_t\clh$. Suppose that $(2)$ holds. Let $\xi \in \clh_{\infty}$. Then, $E_{s}\xi=\xi$ for every $s \in \bbr_{+}$. Since $E_{s}\xi \to 0$ as $s \to \infty$, it follows that $\xi=0$. 
Hence $(2) \implies (3)$. 

 Suppose that $\clh_{\infty}=\{0\}$. Taking orthocomplements, it follows that $\bigcup_{t \geq 0}\ker(V_t)^{*}$ is dense in $\clh$. Since $\{E_t\}_{t \geq 0}$ is norm bounded, it suffices to verify $(2)$ when $\xi \in \bigcup_{s \geq 0}\ker(V_s)^{*}$. 
 Let $s>0$ and let $\xi \in Ker(V_s^*)$ be given. Note that for $t>s$, $Ran(V_t) \subset Ran(V_s)$. Thus, for $t>s$, $E_t\xi=E_tE_s\xi=0$. Consequently, $E_t\xi \to 0$ as $t \to \infty$. This completes the proof of $(3) \implies (2)$. 
 \hfill $\Box$
 
 
 Let $V:=\{V_t\}_{t \geq 0}$ be an isometric representation of $\bbr_{+}$ on a Hilbert space $\clh$. Set $\clh_{\infty}:=\displaystyle \bigcup_{t \geq 0}V_t\clh$. Note that for any $s>0$, $\displaystyle \clh_{\infty}=\bigcap_{t >s}V_t\clh$. More generally, if $F$ is a cofinal subset of $[0,\infty)$, then $\clh_{\infty}=\displaystyle \bigcap_{t \in F}V_t\clh$. It is clear that for $t \geq 0$, $V_t\clh_{\infty}=\clh_{\infty}$ and $V_{t}^{*}\clh_\infty=\clh_\infty$. Thus, $\clh_{\infty}$ is invariant under $\{V_t,V_{t}^{*}: t \geq 0\}$. Also, $\clh_{\infty}^{\perp}$ is invariant under $\{V_t,V_{t}^{*}:t \geq 0\}$. 
 
 
 For $t \geq 0$, let $U_t$ be the restriction of $V_t$ to $\clh_\infty$ and let $W_t$ be the restriction of $V_t$ to $\clh_{\infty}^{\perp}$. If we decompose $\clh$ as $\clh=\clh_{\infty}^{\perp}\oplus \clh_{\infty}$, then $V_t$ is of the form 
\[V_t= \begin{bmatrix}
           W_t & 0 \\
           0 & U_t
           \end{bmatrix}.\] By construction, $U:=\{U_t\}_{t \geq 0}$ is a unitary representation of $\bbr_{+}$ on $\clh_{\infty}$. 
\begin{xrcs}
With the foregoing notation, prove that  $W:=\{W_t\}_{t \geq 0}$ is a pure isometric representation of $\bbr_{+}$ on $\clh_{\infty}^{\perp}$. 
 \end{xrcs} 
 
 Thus, we have proved the following `baby version' of Wold decomposition.
\begin{ppsn}
\label{Baby Wold}
Any isometric representation of $\bbr_{+}$ is unitarily equivalent to a direct sum of a pure isometric representation  of $\bbr_{+}$ and a unitary representation of $\bbr_{+}$. 
\end{ppsn}


Thus, to understand $1$-parameter semigroups of isometries, we need to understand pure isometric representations and unitary representations. We have already talked about  
 unitary representations of $\bbr$. They are in bijective correspondence with non-degenerate representations of $C^{*}(\bbr) \cong C_0(\bbr)$. What can we say about pure semigroups of isometries ? 
 The pleasant feature is that there are only countably many of them and we can list them. This is the content of Cooper's theorem. 

For $t \geq 0$, let $S_t$ be the isometry on $L^{2}((0,\infty))$ defined by \begin{equation}
\label{isometries}
S_{t}(\xi)(s):=\begin{cases}
 \xi(s-t)  & \mbox{ if
} s  \geq t,\cr
   &\cr
    0 &  \mbox{ if } s<t
         \end{cases}
\end{equation}
for $\xi \in L^{2}((0,\infty))$. Then, $S:=\{S_t\}_{t \geq 0}$ is a strongly continuous, pure semigroup of isometries, usually called the shift semigroup on $L^{2}((0,\infty))$. The main result of this section 
is the following theorem. 
\begin{thm}[Cooper]
\label{Cooper}
Let $V:=\{V_t\}_{t \geq 0}$ be a strongly continuous semigroup of isometries on a Hilbert space $\clh$. Suppose that $V$ is pure, i.e. $V_t^{*} \to 0$ strongly as $t \to \infty$. Then, 
there exists a Hilbert space $\mathcal{L}$ and a unitary $U:\clh \to L^{2}((0,\infty)) \otimes \mathcal{L}$ such that 
 for $t \geq 0$, 
$
UV_tU^{*}=S_t \otimes 1$.
\end{thm}

The usual proofs available in the literature make essential use of unbounded operators. The proof that is provided here is more algebraic.  The idea behind the proof is to look for  an algebra whose non-degenerate
representations are in $1$-$1$ correspondence with pure semigroups of isometries. It turns out that the algebra that we seek is $C_0(\bbr) \rtimes \bbr$, where the action of $\bbr$ is by translation.
Thanks to the abstract version of the Stone-von Neumann theorem (Thm. \ref{Stone abstract}), $C_0(\bbr) \rtimes \bbr \cong \mathcal{K}(L^{2}(\bbr))$ and the latter algebra has only countably
many non-degenerate representations up to unitary equivalence. We make this precise in what follows. 


\textbf{Unitary dilations:} Let $V:=\{V_t\}_{t \geq 0}$ be an isometric representation of $\bbr_+$ on a Hilbert space $\clh$. Suppose $U:=\{U_t\}_{t \in \bbr}$ is a unitary representation on a Hilbert space $\mathcal{K}$ and suppose $j:\clh \to \mathcal{K}$ is an isometry. We say that the triple $(U,\mathcal{K},j)$ is a unitary dilation of $V$ if for $t \geq 0$ and $\xi \in \clh$, 
\[
U_t(j(\xi))=j(V_t\xi).\]
We usually suppress the notation $j$ and regard $\clh$ as a closed subspace of $\mathcal{K}$ and by abusing notation, we write the above equality as $V_t\xi=U_t\xi$ whenever $t \geq 0$ and $\xi \in \clh$. Let $(U,\mathcal{K})$ be a unitary dilation of $V$. We say that $(U,\mathcal{K})$ is \emph{minimal} if $\bigcup_{t \geq 0}U_{-t}\clh$ is dense in $\mathcal{K}$.  

Let $(U^{(1)},\clk_1)$ and $(U^{(2)},\clk_2)$ be two unitary dilations of $V$. We say that they are unitarily equivalent if there exists a unitary $X:\clk_1 \to \clk_2$ such that $X\clh=\clh$ and for $t \in \bbr$, $XU^{(1)}_tX^{*}=U^{(2)}_{t}$. 

\begin{thm}
Let $V:=\{V_t\}_{t \geq 0}$ be an isometric representation of $\bbr_{+}$ on a Hilbert space $\clh$. Then, $V$ admits a minimal unitary dilation. Moreover, any two minimal unitary dilations of $V$ are unitarily equivalent. 
\end{thm}
\textit{Proof.} First, we prove the uniqueness. Let $(U^{(1)},\clk_1)$ and $(U^{(2)},\clk_2)$ be two minimal unitary dilations of $V$. Let $\mathcal{D}_{i}:=\bigcup_{t \geq 0}U^{(i)}_{-t}\clh$. 
Then, $\mathcal{D}_i$ is a dense subspace of $\clk_i$. If we make use of Remark \ref{KRP}, we see that there exists a unique unitary $X:\clk_1 \to \clk_2$ such that 
\[
XU^{(1)}_{-t}\xi=U^{(2)}_{-t}\xi\]
for $t \geq 0$ and $\xi \in \clh$. It is routine to verify that $X$ implements the required equivalence between $(U^{(1)},\clk_1)$ and $(U^{(2)},\clk_2)$. 

Next, we construct a dilation that is minimal. Define an equivalence relation on $\clh \times [0,\infty)$ as follows. We say $(\xi,s) \sim (\eta,t)$ if $V_t\xi=V_s\eta$. Denote the set of equivalence classes by $\mathcal{K}_{\infty}$. Then, $\mathcal{K}_{\infty}$ has the structure of an inner product space where addition, scalar multiplication and inner product are given by 
\begin{align*}
[(\xi,s)]+[(\eta,t)]&=[(V_t\xi+V_s\eta,s+t)]\\
\lambda[(\xi,s)]&=[(\lambda\xi,s)]\\
\langle [(\xi,s)]|[(\eta,t)]\rangle&=\langle V_t\xi|V_s\eta \rangle.
\end{align*}
Denote the completion of $\clk_{\infty}$ by $\clk$. Let $j:\clh \to \clk$ be defined by $j(\xi):=[(\xi,0)]$. Then, $j$ is an isometry. Through $j$, we view $\clh$ as a closed subspace of $\clk$. 

For $t \geq 0$, let $U_t:\clk_\infty \to \clk_\infty$ be defined by 
\[
U_t([(\xi,s)]=[(V_t\xi,s)].\]
It is clear that $U_t$ is well defined, linear and preserves the inner product. Clearly, for $\xi \in \clh$, $U_t\xi=V_t\xi$. More precisely, $U_t(j(\xi))=j(V_t\xi)$ for $t \geq 0$ and $\xi \in \clh$.

We claim that for $t \geq 0$,  $U_t$ is onto. Let $[(\eta,s)] \in \clk_\infty$ be given. Then, by definition, 
\begin{equation}
\label{formula for inverse}
[(\eta,s)]=[(V_t\eta,s+t)]=U_t([(\eta,s+t)]).\end{equation}
The above equality implies that $U_t$ is onto. %Moreover, $U_{t}^{-1}([(\eta,s)]=[(\eta,s+t)]$. 
Thus, $U_t$ extends to a unitary operator on $\clk$ which we denote again by $U_t$. Clearly, $U_sU_t=U_{s+t}$ for $s,t \geq 0$. Thanks to Eq. \ref{formula for inverse}, we have $\bigcup_{t \geq 0}U_{t}^{-1}\clh=\bigcup_{t \geq 0}\{[(\xi,t):\xi \in \clh\}=\clk_\infty$ and the latter set is dense in $\clk$ by definition. 

Now, we check the strong continuity of the map $[0,\infty) \ni t \to U_t \in B(\clh)$. It suffices to prove that for every $\xi \in \clh$ and $s \geq 0$, the map $[0,\infty) \ni t \to U_t([(\xi,s)]) \in \clk$ is continuous. This follows from the strong continuity of $V$ and the inequality 
\[
|| U_{t_1}([(\xi,s)])-U_{t_2}([(\xi,s)])||=||V_{t_1}\xi-V_{t_2}\xi||.\]
Thus, $U:=\{U_t\}_{t \geq 0}$ is a strongly continuous group of unitaries. Extend $U$ to the whole of $\bbr$ and denote the extension again by $U$. Then, $U$ is clearly `the minimal unitary dilation' of $V$. 
\hfill $\Box$

\begin{xrcs}
Let $S:=\{S_t\}_{t \geq 0}$ be the shift semigroup on $L^{2}((0,\infty))$. Suppose $\lambda:=\{\lambda_t\}_{t \in \bbr}$ is the regular representation of $\bbr$ on $L^{2}(\bbr)$. Prove that $\lambda$ is the minimal unitary dilation of $S$. 
\end{xrcs}
\textit{Hint:} View $L^{2}((0,\infty))$ as a subspace of $L^{2}(\bbr)$ in the obvious way. It suffices to check that $\lambda:=\{\lambda_t\}_{t \in \bbr}$ has `the defining properties' for it to be the minimal unitary dilation. 

We will proceed towards the proof of Cooper's theorem. Let $V:=\{V_t\}_{t \geq 0}$ be a strongly continuous semigroup of isometries on a Hilbert space $\clh$. Assume that $V$ is pure. The isometric representation $V$ will be fixed for the rest of this section. 
Let $(U,\clk)$ be the minimal unitary dilation of $V$. For $t \in \bbr$, let $E_t$ be the orthogonal projection onto $U_t\clh$. 

Observe the following properties regarding the family $\{E_t:t \in \bbr\}$ of projections. 
\begin{enumerate}
\item[(i)] Since $U_t\clh \subset \clh$ for every $ t\geq 0$, it follows that if $s\leq t$, then $U_t\clh \subset U_s\clh$. This means that if $s \leq t$, then $E_s \geq E_t$. In particular, the family $\{E_t: t \in \bbr\}$ is a commuting family of projections. 
\item[(ii)] For $s,t \in \bbr$, $U_{s+t}\clh=U_s(U_t\clh)$. Thus, we have the following covariance relation. For $s \in \bbr$ and $t \in \bbr$, 
\[
U_sE_tU_{s}^{*}=E_{s+t}.\]
\item[(iii)] The map $\bbr \ni t \to E_{t} \in B(\clk)$ is strongly continuous. This follows from the strong continuity of $\{U_t\}_{t \in \bbr}$ and from the fact that $E_t=U_tE_0U_{t}^{*}$. 
\end{enumerate}

The key ingredient in the proof of Cooper's theorem is to understand the algebra generated by $\{E_t:t \in \bbr\}$. However, if we take the $C^{*}$-algebra generated by the family $\{E_t: t \in \bbr\}$, then it is not separable, and in particular its spectrum
will be `too big'. Taking inspiration from the study of group $C^{*}$-algebra, we look at the $C^{*}$-algebra generated by $\{\int f(t)E_t dt: f \in C_{c}(\bbr)\}$. 

Let us introduce a few notation. Adjoint $\infty$ at the right end of $\bbr$ and endow $(-\infty,\infty]$ with the usual topology. For $\phi \in C_0((-\infty,\infty])$ and $s \in \bbr$, let $R_s(\phi) \in C_{0}((-\infty,\infty])$ be defined by 
\[
R_s\phi(x)=\phi(x-s).\]
As usual, the convention is that $\infty\pm t=\infty$ for all $t \in \bbr$. Then, $R:=\{R_s\}_{s \in \bbr}$ defines an action of $\bbr$ on $C_0((-\infty,\infty])$. 


For $f \in C_{c}(\bbr)$, let $\widetilde{f}:(-\infty,\infty] \to \bbc$ be defined by 
\[
\widetilde{f}(x):=\int_{-\infty}^{x}f(t)dt.\]
Then, $\widetilde{f} \in C_0((-\infty,\infty])$ for every $f \in C_{c}(\bbr)$. 

Let $\mathcal{D}$ be the $C^{*}$-algebra generated by $\Big \{\int f(t)E_t dt: f \in C_{c}(\bbr)\Big\}$. Note that $\mathcal{D}$ is commutative. %The heart of the proof of Thm. \ref{Cooper} lies in the following two propositions. 

\begin{lmma}
\label{description of characters}
Let $\chi$ be a character of the commutative $C^{*}$-algebra $\mathcal{D}$. Then, there exists a unique $x=:x_{\chi} \in (-\infty,\infty]$ such that 
\[
\chi\Big(\int f(t)E_tdt\Big)=\int_{-\infty}^{x}f(t)dt\]
for all $f \in C_{c}(\bbr)$. 
\end{lmma}
\textit{Proof.} Let $\chi$ be a character of $\mathcal{D}$. Note that the map 
\[
C_{c}(\bbr) \ni f \to \chi\Big(\int f(t)E_tdt\Big) \in \bbc\]
admits an extension to $L^{1}(\bbr)$. Therefore, there exists $\phi \in L^{\infty}(\bbr)$ such that 
\[
\chi\Big(\int f(t)E_tdt\Big)=\int f(t)\phi(t)dt\]
for $f \in C_{c}(\bbr)$. 

Since $\chi$ is non-zero, there exists real numbers $a$ and $b$ with $a<b$ and a function $g \in C_{c}(\bbr)$ such that $supp(g) \subset (a,b)$ and $\int g(s)\phi(s)dt=1$. 
We claim that $\phi(t)=1$ for almost all $t \in (-\infty,a)$. 

Let $f \in C_{c}(\bbr)$ be any function such that $supp(f) \subset (-\infty,a)$. 
Note that $E_tE_s=E_s$ whenever $t<a<s$. Therefore, 
\begin{align*}
\Big(\int f(t)E_tdt\Big)\Big(\int g(s)E_sds\Big)&=\Big(\int_{-\infty}^{a}f(t)E_tdt\Big)\Big(\int_{a}^{b}g(s)E_sds\Big)\\
&=\Big(\int_{(t,s) \in (-\infty,a) \times (a,b)} f(t)g(s)E_tE_sdtds\Big)\\
&=\Big(\int_{-\infty}^{a}f(t)dt\Big)\int_{-\infty}^{\infty}g(s)E_sds.
\end{align*}
Applying $\chi$ to the above equation and appealing to the fact that $\chi$ is a character, we deduce that 
\begin{align*}
\int_{-\infty}^{a}f(t)dt& = \Big(\int_{-\infty}^{a}f(t)dt\Big)\Big(\int_{a}^{b} g(s)\phi(s)ds\Big)\\
&=\Big(\int_{-\infty}^{a}f(t)dt\Big)\chi\Big(\int g(s)E_s ds\Big)\\
&=\chi\Big(\Big(\int_{-\infty}^{a}f(t)dt\Big)\Big(\int g(s)E_sds\Big)\Big) \\
&=\chi\Big(\Big(\int f(t)E_tdt \Big)\Big(\int g(s)E_sds\Big)\Big)\\
&=\chi\Big(\int f(t)E_tdt \Big)\chi\Big(\int g(s)E_sds\Big)\\
&=\int_{-\infty}^{a}f(t)\phi(t)dt \int_{a}^{b} g(s)\phi(s)ds \\
&=\int_{-\infty}^{a}f(t)\phi(t)dt. 
\end{align*}
Since $f$ is arbitrary, we can conclude that $\phi(t)=1$ for almost all $t \in (-\infty,a)$.

Define \[A:=\{a \in \bbr: \phi(t)=1 \textrm{~~for almost all $t \in (-\infty,a)$}\}.\] We have shown that $0 \in A$. Set $x:=\sup(A)$. Clearly, $\phi(t)=1$ for almost all $t \in (-\infty,x]$.  If $x=\infty$, then $\phi=1$ a.e. and we are done. 
Suppose that $x<\infty$. 

We claim that $\phi(t)=0$ for almost all $t \in (x,\infty)$. Suppose not. Then, there exist real numbers $a,b$ with $b>a>x$ and a function $g \in C_{c}(\bbr)$ with $supp(g) \subset (a,b)$ and $\int_{a}^{b} g(s)\phi(s)ds=1$. Arguing as before, we 
can conclude that $\phi(t)=1$ for almost all $t \in (-\infty,a)$. Then, $A \ni a >x$ which is a contradiction. This proves the claim. In other words, $\phi(t)=1_{(-\infty,x]}(t)$ a.e. Hence the proof. \hfill $\Box$

\begin{ppsn}
\label{covariant Cooper}
Keep the foregoing notation. 
\begin{enumerate}
\item[(1)] There exists a unique $*$-homomorphism $\pi:C_0((-\infty,\infty]) \to B(\clk)$ such that for $f \in C_{c}(\bbr)$, 
\[
\pi(\widetilde{f})=\int f(t)E_t dt.\]
\item[(2)] The pair $(\pi,U)$ is covariant, i.e. for $\phi \in C_0((-\infty,\infty])$ and $s \in \bbr$, 
\[
U_s\pi(\phi)U_s^{*}=\pi(R_s\phi).\]
\end{enumerate}
\end{ppsn}
\textit{Proof.} Uniqueness of the homomorphism $\pi$ follows from the fact that $\{\widetilde{f}:f \in C_{c}(\bbr)\}$ generates $C_0((-\infty,\infty])$. This can be proved using Stone-Weierstrass theorem. For the existence, let us denote the map 
\[
\widehat{D} \ni \chi \to x_\chi \in (-\infty,\infty]\]
by $T$. Note that the map $T$ is continuous. For $\phi \in C_0((-\infty,\infty])$, set \[\pi(\phi):=G(\phi \circ T).\] Here, $G: C(\widehat{\mathcal{D}}) \to \mathcal{D} \subset B(\clk)$ is the Gelfand transformation. It follows from definitions that for $f \in C_{c}(\bbr)$, 
\[
\pi(\widetilde{f})=\int f(t)E_t dt.\]

Let $s \in \bbr$ and $f \in C_{c}(\bbr)$ be given. Note that $\widetilde{R_sf}=R_s(\widetilde{f})$. Calculate as follows to observe that 
%\iffalse
\begin{align*}
U_s\pi(\widetilde{f})U_s^{*}&=U_{s} \Big (\int f(t)E_tdt\Big)U_{s}^{*}\\
&=\int f(t)U_sE_tU_s^{*}dt\\
&=\int f(t)E_{s+t}dt\\
&=\int f(t-s)E_tdt\\
&=\int R_s(f)(t)E_tdt\\
&=\pi(\widetilde{R_s(f)})=\pi(R_s(\widetilde{f})).
\end{align*}
Since $\{\widetilde{f}:f \in C_{c}(\bbr)\}$ generates $C_0((-\infty,\infty])$, it follows that $U_s\pi(\phi)U_{s}^{*}=\pi(R_s\phi)$ for $s \in \bbr$ and $\phi \in C_0((-\infty,\infty])$. This completes the proof. \hfill $\Box$
%\fi

There are a couple of crucial observations that we need to make. Recall that we have assumed $V$ is pure. 
\begin{lmma}
\label{Cooper vanishing at infinity}
For every $\xi \in \clk$, $E_t\xi \to 0$ as $t \to +\infty$.
\end{lmma}
\textit{Proof.} Note that for $s \geq 0$, $V_s=E_0U_sE_0$ and $V_{s}^{*}=E_0U_s^*E_0$. 
As $\{E_t:t \in \bbr\}$ is norm bounded,  It suffices to verify the statement when $\xi \in \bigcup_{s \geq 0}U_{s}^{*}\clh$. Let $s \geq 0$ and $\eta \in \clh$ be given. Note that for large $t$.
\[
E_{t}U_{s}^{-1}\xi=U_{t}E_0U_{t-s}^{*}E_0\xi=U_{t}V_{t-s}^{*}\xi.\]
Hence, for large $t$, 
\[
||E_{t}U_{s}^{-1}\xi|| \leq ||V_{t-s}^{*}\xi||.\]
The purity of $V$ implies that $V_{t-s}^{*}\xi \to 0$ as $t \to \infty$. Hence the proof. \hfill $\Box$


\begin{lmma}
\label{non-degenerate Cooper}
Let $\pi$ be the representation of $C_0((-\infty,\infty])$ described in Thm. \ref{covariant Cooper}. Then, $\pi$ restricted to $C_{0}(\bbr)$ is non-degenerate. 
\end{lmma}
\textit{Proof.} Suppose not. Then, there exists a non-zero vector $\eta$ such that 
\[
\langle \pi(\phi)\xi|\eta\rangle=0\]
for every $\phi \in C_{c}(\bbr)$ and $\xi \in \clk$. 

Let $\xi \in \clk$ be arbitrary. Let $\omega_\xi:C_{c}(\bbr) \to \bbc$ be defined by 
\[
\omega_\xi(f)=\int f(t)\langle E_t\xi|\eta \rangle dt.\]
Define $I:C_{c}(\bbr) \to \bbc$ by \[I(f)=\int f(t)dt.\] Suppose $f \in \ker(I)$. Then, $\widetilde{f} \in C_{c}(\bbr)$. The fact that $\langle \pi(\widetilde{f})\xi|\eta \rangle=0$ implies that 
$\int f(t)\langle \xi|E_t\eta\rangle=0$, i.e $\omega_\xi(f)=0$. Consequently, $\ker(I)=\ker(\omega_\xi)$. Thus, there exists $c_{\xi} \in \bbc$ such that for all $f \in C_{c}(\bbr)$, 
\[
\int f(t)\langle E_t\xi|\eta \rangle dt=c_{\xi}\int f(t)dt.\]
As the above equality happens for every $f \in C_{c}(\bbr)$ and the map $\bbr \to t \to \langle E_t\xi|\eta \rangle \in \bbc$ is continuous, we conclude that $\langle E_t\xi|\eta\rangle =c_{\xi}$ for every $t \in \bbr$. 
Lemma \ref{Cooper vanishing at infinity} implies that $\langle E_t\xi|\eta \rangle=0$ for every $t \in \bbr$. 
 
 Thus, $\eta$ is orthogonal to $U_t\clh$ for every $t \in \bbr$. However, $\bigcup_{t \geq 0}U_{-t}\clh$ is dense in $\clk$. Thus, $\eta$ is orthogonal to every vector of $\clk$ forcing $\eta=0$. This contradiction
 completes the proof. \hfill $\Box$
 
 The final step that is required is the following lemma. Define
 \[
 \pi(C_{c}((0,\infty)))\clk:=\overline{span\{\pi(\phi)\xi:\phi \in C_{c}((0,\infty)), \xi \in \clk\}}.\]
 
 \begin{lmma}
 \label{projection onto the given space}
 The orthogonal projection onto $\pi(C_{c}((0,\infty)))\clk$ is $E_0$, or in other words, $\clh=\pi(C_{c}((0,\infty)))\clk$.  
  \end{lmma}
  \textit{Proof.} Let $C_{c}^{\infty}((0,\infty))$ be the algebra of smooth functions $f$ on $\bbr$ that vanishes outside a compact subset contained in $(0,\infty)$. Clearly, $\pi(C_{c}^{\infty}((0,\infty)))\clk=\pi(C_{c}((0,\infty)))\clk$. 
 Let $\phi \in C_{c}^{\infty}((0,\infty))$ be given. Observe that 
 \[
 E_0\pi(\phi)=E_0\Big(\int_{0}^{\infty}\phi^{'}(t)E_tdt\Big)=\int_{0}^{\infty} \phi^{'}(t)E_0E_tdt=\int \phi^{'}(t)E_tdt=\pi(\phi).\]
 The above equality implies that $\pi(C_{c}^{\infty}((0,\infty)))\clk \subset \clh$. 
 
 Suppose that the inclusion $\pi(C_c^{\infty}((0,\infty)))\clk \subset \clh$ is proper. Then, there exists a non-zero vector $\eta \in \clh$ such that for every $\phi \in C_{c}^{\infty}((0,\infty))$ and $\xi \in \clk$,
 \[
 \int_{0}^{\infty} \phi^{'}(t)\langle E_t\xi|\eta \rangle dt=0.\]
 Let $f \in C_{c}^{\infty}((0,\infty))$ be such that $\int f(t)dt=0$. Define $\phi(x):=\int_{-\infty}^{x}f(t)dt$. Then, $\phi^{'}(x)=f(x)$. The hypothesis implies that 
 \[
 \int f(t)\langle E_t\xi|\eta \rangle dt=0.\]
 Thus, arguing as in Lemma \ref{non-degenerate Cooper}, we can conclude that $\langle E_t\xi|\eta\rangle =0$ for every $t \in (0,\infty)$. Therefore, $E_t\eta=0$ for every $t>0$. Since $E_t\eta \to E_0\eta=\eta$ as $t \to 0$, we conclude that $\eta=0$, which is a contradiction. Hence the proof. \hfill $\Box$
 
 
 \textit{Proof of Thm. \ref{Cooper}:} Let $(\pi,U)$ be the covariant representation of the dynamical system $(C_0((-\infty,\infty]),\bbr,R)$ constructed in Prop. \ref{covariant Cooper}. Thanks to Lemma \ref{non-degenerate Cooper}, $\pi$ restricted to $C_0(\bbr)$ is non-degenerate. Thus, we obtain a non-degenerate covariant representation $(\pi,U)$ of the dynamical system $(C_0(\bbr),\bbr,R)$ on $\clk$. By the Stone-von Neumann theorem (Ex. \ref{Stone another}), there exists a Hilbert space $\mathcal{L}$ and a unitary $X:\clk \to L^{2}(\bbr) \otimes \mathcal{L}$ such that for $\phi \in C_0(\bbr)$ and $s \in \bbr$, 
 \begin{align*}
 X\pi(\phi)X^{*}&=M(\phi)\otimes 1\\
 XU_sX^{*}&=\lambda_s \otimes 1.
 \end{align*}
 Here, $M:C_0(\bbr) \to B(L^{2}(\bbr))$ is the multiplication representation and $\lambda:=\{\lambda_s\}_{s \in \bbr}$ is the regular representation of $\bbr$ on $L^{2}(\bbr)$. It is clear from the above equation and Lemma \ref{projection onto the given space} that $X$ takes $\clh$ onto $L^{2}((0,\infty))\otimes \mathcal{L}$. Let $Y$ be the restriction of $X$ to $\clh$. Then, $Y$ is a unitary between $\clh$ and $L^{2}((0,\infty))\otimes \mathcal{L}$. 
 
 Let $S:=\{S_t\}_{t \geq 0}$ be the shift semigroup on $L^{2}((0,\infty))$. 
 For $t \geq 0$, $S_t \otimes 1$ is the restriction of $\lambda_t \otimes 1$ to $L^{2}((0,\infty))\otimes \mathcal{L}$. As $U:=\{U_t\}_{t \geq 0}$ is a dilation of $V:=\{V_t\}_{t \geq 0}$, for $t \geq 0$, $V_t$ is the restriction of $U_t$ to $\clh$. As $X$ intertwines $\{U_t\}_{t \geq 0}$ and $\{\lambda_t\otimes 1\}_{t \geq 0}$, it follows that $Y$ intertwines $\{V_t\}_{t \geq 0}$ and $\{S_t \otimes 1\}_{t \geq 0}$. The proof is complete. \hfill $\Box$
 
 With Cooper's theorem in hand, the Wold decomposition can now be restated as follows. 
 
 \begin{thm}[Wold-decomposition]
 Let $V:=\{V_{t}\}_{t \geq 0}$ be a strongly continuous semigroup of isometries on a Hilbert space $\clh$. Then, there exist Hilbert spaces $\mathcal{L}$ and $\clk$, a  unitary representation $U:=\{U_t\}_{t \geq 0}$ on $\clk$ and a unitary $X:\clh \to \big(L^{2}((0,\infty))\otimes \mathcal{L} \big)\oplus \mathcal{K}$ such that for every $t \geq 0$,
 \[XV_tX^{*}=\begin{bmatrix}
 S_t \otimes 1 & 0 \\
 0 & U_t
 \end{bmatrix}.\] 
 
 
 
 \end{thm}
 
 
 
 % \newpage
  